% GENERATED FILE. Edit canonical lessons, then run npm run generate:books.
% canonical-source-hash: fnv1a64:3fababac8ba73404
% canonical-lessons: RU-C236-schitat, RU-C236-izmeryat, RU-C236-sravnivat, RU-C236-proveryat, RU-C236-zapolnyat

\chapter{To count, To measure, To compare, To check, To fill in}
\label{ch:ru-to-count-to-measure-to-compare-to-check-to-fill-in}
\hlchaptermodality{236}

\begin{chapteropening}
\textbf{By the end of this chapter:} \emph{I can say to count, to measure, to compare, to check and to fill in.}

Complete the last lesson of chapter 236: I can say to count, to measure, to compare, to check and to fill in.
\end{chapteropening}

\section[\texorpdfstring{\ru{считать} (schitát)}{считать (schitát)}]{\ru{считать} (schitát) — to count; to consider}

\label{lesson:RU-C236-schitat}



\begin{quote}
Before the new one: say the Russian for to smile, then the Russian for to be angry.
\end{quote}

\subsection*{You'll want to know: \ru{считать}}
\textbf{\ru{считать}} (\emph{schitát}) — \textquotedblleft{}to count; to consider\textquotedblright{}.

\begin{cousinweb}
\textbf{\ru{считать}} is \emph{\ru{с}-} with \emph{\ru{читать}}, \textquotedblleft{}to read\textquotedblright{}, which once meant \textquotedblleft{}to count\textquotedblright{}, as in \emph{\ru{число}}, \textquotedblleft{}a number\textquotedblright{}. Vasmer links \emph{\ru{читать}} with an old Slavic verb meaning \textquotedblleft{}to count, to honour\textquotedblright{}.
\end{cousinweb}

\begin{grammarlens}[title={What you've built}]
One more word for this chapter.
\end{grammarlens}

\subsection*{Your turn}
\begin{itemize}
\raggedright
  \item \emph{Say it:} \emph{schitát}
  \item \emph{Say it:} \emph{schitát}, once more
  \item \emph{Say it:} say \emph{schitát}
  \item \emph{Recall:} say the Russian for to smile, then the Russian for to be angry, then say \emph{schitát} again
\end{itemize}

\begin{tcolorbox}[breakable,colback=teal!4,colframe=teal!35!black,title={Before you move on}]
What does \ru{считать} mean? (\textquotedblleft{}To count; to consider\textquotedblright{}.) Say it once more.
\end{tcolorbox}
\section[\texorpdfstring{\ru{измерять} (izmeryát)}{измерять (izmeryát)}]{\ru{измерять} (izmeryát) — to measure}

\label{lesson:RU-C236-izmeryat}



\begin{quote}
Before the new one: say the Russian for to be angry, then the Russian for to count; to consider.
\end{quote}

\subsection*{You'll want to know: \ru{измерять}}
\textbf{\ru{измерять}} (\emph{izmeryát}) — \textquotedblleft{}to measure\textquotedblright{}.

\begin{cousinweb}
\textbf{\ru{измерять}} is \emph{\ru{из}-}, \textquotedblleft{}out, through\textquotedblright{}, with \emph{\ru{мерять}}, from \emph{\ru{мера}}, \textquotedblleft{}a measure\textquotedblright{}, which Vasmer relates to Latin \emph{mētior}, \textquotedblleft{}I measure\textquotedblright{}.
\end{cousinweb}

\begin{grammarlens}[title={What you've built}]
One more word for this chapter.
\end{grammarlens}

\subsection*{Your turn}
\begin{itemize}
\raggedright
  \item \emph{Say it:} \emph{izmeryát}
  \item \emph{Say it:} \emph{izmeryát}, once more
  \item \emph{Say it:} say \emph{izmeryát}
  \item \emph{Recall:} say the Russian for to be angry, then the Russian for to count; to consider, then say \emph{izmeryát} again
\end{itemize}

\begin{tcolorbox}[breakable,colback=teal!4,colframe=teal!35!black,title={Before you move on}]
What does \ru{измерять} mean? (\textquotedblleft{}To measure\textquotedblright{}.) Say it once more.
\end{tcolorbox}
\section[\texorpdfstring{\ru{сравнивать} (srávnivat)}{сравнивать (srávnivat)}]{\ru{сравнивать} (srávnivat) — to compare}

\label{lesson:RU-C236-sravnivat}



\begin{quote}
Before the new one: say the Russian for to count; to consider, then the Russian for to measure.
\end{quote}

\subsection*{You'll want to know: \ru{сравнивать}}
\textbf{\ru{сравнивать}} (\emph{srávnivat}) — \textquotedblleft{}to compare\textquotedblright{}.

\begin{cousinweb}
\textbf{\ru{сравнивать}} is \emph{\ru{с}-} with \emph{-\ru{равнивать}}, from \emph{\ru{равный}}, \textquotedblleft{}equal\textquotedblright{}, the Church Slavonic shape of \emph{\ru{ровный}}, \textquotedblleft{}level\textquotedblright{}: to compare is to set things level.
\end{cousinweb}

\begin{grammarlens}[title={What you've built}]
One more word for this chapter.
\end{grammarlens}

\subsection*{Your turn}
\begin{itemize}
\raggedright
  \item \emph{Say it:} \emph{srávnivat}
  \item \emph{Say it:} \emph{srávnivat}, once more
  \item \emph{Say it:} say \emph{srávnivat}
  \item \emph{Recall:} say the Russian for to count; to consider, then the Russian for to measure, then say \emph{srávnivat} again
\end{itemize}

\begin{tcolorbox}[breakable,colback=teal!4,colframe=teal!35!black,title={Before you move on}]
What does \ru{сравнивать} mean? (\textquotedblleft{}To compare\textquotedblright{}.) Say it once more.
\end{tcolorbox}
\section[\texorpdfstring{\ru{проверять} (proveryát)}{проверять (proveryát)}]{\ru{проверять} (proveryát) — to check}

\label{lesson:RU-C236-proveryat}



\begin{quote}
Before the new one: say the Russian for to measure, then the Russian for to compare.
\end{quote}

\subsection*{You'll want to know: \ru{проверять}}
\textbf{\ru{проверять}} (\emph{proveryát}) — \textquotedblleft{}to check\textquotedblright{}.

\begin{cousinweb}
\textbf{\ru{проверять}} is \emph{\ru{про}-}, \textquotedblleft{}through\textquotedblright{}, with \emph{-\ru{верять}}, from \emph{\ru{вера}}, \textquotedblleft{}trust, faith\textquotedblright{}: to check is to test whether a thing can be trusted.
\end{cousinweb}

\begin{grammarlens}[title={What you've built}]
One more word for this chapter.
\end{grammarlens}

\subsection*{Your turn}
\begin{itemize}
\raggedright
  \item \emph{Say it:} \emph{proveryát}
  \item \emph{Say it:} \emph{proveryát}, once more
  \item \emph{Say it:} say \emph{proveryát}
  \item \emph{Recall:} say the Russian for to measure, then the Russian for to compare, then say \emph{proveryát} again
\end{itemize}

\begin{tcolorbox}[breakable,colback=teal!4,colframe=teal!35!black,title={Before you move on}]
What does \ru{проверять} mean? (\textquotedblleft{}To check\textquotedblright{}.) Say it once more.
\end{tcolorbox}
\section[\texorpdfstring{\ru{заполнять} (zapolnyát)}{заполнять (zapolnyát)}]{\ru{заполнять} (zapolnyát) — to fill in}

\label{lesson:RU-C236-zapolnyat}



\begin{quote}
Before the new one: say the Russian for to compare, then the Russian for to check.
\end{quote}

\subsection*{You'll want to know: \ru{заполнять}}
\textbf{\ru{заполнять}} (\emph{zapolnyát}) — \textquotedblleft{}to fill in\textquotedblright{}.

\begin{cousinweb}
\textbf{\ru{заполнять}} is \emph{\ru{за}-} with \emph{-\ru{полнять}}, from \emph{\ru{полный}}, \textquotedblleft{}full\textquotedblright{}: to fill in.
\end{cousinweb}

\begin{grammarlens}[title={What you've built}]
One more word for this chapter.
\end{grammarlens}

\subsection*{Your turn}
\begin{itemize}
\raggedright
  \item \emph{Say it:} \emph{zapolnyát}
  \item \emph{Say it:} \emph{zapolnyát}, once more
  \item \emph{Say it:} say \emph{zapolnyát}
  \item \emph{Recall:} say the Russian for to compare, then the Russian for to check, then say \emph{zapolnyát} again
\end{itemize}

\begin{tcolorbox}[breakable,colback=teal!4,colframe=teal!35!black,title={Before you move on}]
What does \ru{заполнять} mean? (\textquotedblleft{}To fill in\textquotedblright{}.) Say it once more.
\end{tcolorbox}
